\section*{Author Section 6.2: first-approximate factorization of $M_c^R$}
Use Section5.3's terminology and extend the recursive approximate factorization.
\paragraph{Definition 6.11.} Fix $R\in\binom{[n]}{\le d/2}$. For $i=0,\ldots,\tau$, define the left-$i$-factor $L_{R,i}$, of dimension $\binom{[n]}{d/2}\times\binom{[n]}{\le d/2}$, by
\[L_{R,i}(I,A)=\begin{cases}0,&R\not\subseteq I\cap A,\\\displaystyle\sum_{\substack{T:|V(T)\cup I\cup A|\le\tau,\ A=S_l(I,A;T)\\T\cap E(A)=\emptyset,\ (I,A;T)\text{ left-generated}\\e_{I,A}(T)=i}}(\omega/n)^i\chi_T,&\text{o.w.}\end{cases}\tag{6.20}\]
$(L_{R,j})^\top$ is the right-$j$-factor. $\widetilde L_R=(L_{R,0},\ldots,L_{R,\tau})$ is the left factor and $(\widetilde L_R)^\top$ the right factor. These do not depend on $c$.
\paragraph{Definition 6.12.} Let $D^\tau$ be the constant diagonal matrix
\[\operatorname{diag}((\omega/n)^{|A|/2})_{A:|A|\le d/2}\otimes\mathrm{Id}_{\{0,\ldots,\tau\}\times\{0,\ldots,\tau\}},\]
of dimension $(\binom{[n]}{\le d/2}\times(\tau+1))\times(\binom{[n]}{\le d/2}\times(\tau+1))$.
\paragraph{Definition 6.13 (Goal factorization).} Find a middle matrix $Q_c^R$ of that same dimension such that approximately
\[M_c^R\approx\widetilde L_R[D^\tau Q_c^R D^\tau](\widetilde L_R)^\top.\tag{6.21}\]
\paragraph{Remark 6.14.} Unlike Section5.3, further distinguish a parameter pair in $\{0,\ldots,\tau\}^2$: in(6.13), and any exact pseudo-expectation generated by Section3.2, $a=|V(T)\cup I\cup J|$ appears nestedly. Previously $(\omega/n)^a=(\omega/n)^{e(R_l)+|V(R_m)|+e(R_r)}$ naturally factors left/middle/right. Here terms like $\binom{a+l-d}c\binom{n-a}{l-c}$ are not log-additive in $a$. The chosen generating function makes $\mathbb E[Q_{0,0}^R]$ positive in this harder situation, as Proposition6.28/Corollary6.30 will show. Separate the main factor $(\omega/n)^a$ as before, and leave $Y_c(r,a)$ to the middle matrix indexed by $(e_l,e_r)$.
\paragraph{Definition 6.15 (First-approximate factorization).} $Q_{c,0}^R$ is a $(\tau+1)\times(\tau+1)$ block matrix, each block of dimension $\binom{[n]}{\le d/2}\times\binom{[n]}{\le d/2}$, zero outside the principal minor $\mathcal S^R\times\mathcal S^R$ with
\[\mathcal S^R:=\{(A,i)\in\binom{[n]}{\le d/2}\times\{0,\ldots,\tau\}:A\supseteq R,\ |A|+i\ge d/2\}.\tag{6.22}\]
Inside this minor,
\begin{align*}Q_{c,0}^R((A,i),(B,j))=\sum_{\substack{T_m:|V(T_m)\cup A\cup B|\le\tau\\A,B\in\mathrm{mSep}_{A,B}(T_m)}}&(\omega/n)^{|V(T_m)\cup A\cup B|-(|A|+|B|)/2}\\&\cdot Y_c(|R|,|V(T_m)\cup A\cup B|+(i+j))\chi_{T_m}.\tag{6.23}\end{align*}
Define $\widetilde L_R[D^\tau Q_{c,0}^R D^\tau](\widetilde L_R)^\top$ as the first approximate factorization of $M_c^R$.
\paragraph{Remark 6.16.} (1) $\mathcal S^R$ is independent of $c$; $|A|+i\ge d/2$ is natural since $i$ is intended as $|V(T')\setminus A|\ge|I|-|A|$ for a left ribbon. Otherwise that column in $\widetilde L_R$ is zero. Similarly for $j$. (2) $Q_{c,0}^R$ is supported where $|A|=|B|$. (3) For an outer-canonical product in the above matrix product, Remark5.12 gives $|V(T_m)\cup A\cup B|+(i+j)=|V(T)\cup I\cup J|$ for the canonical decomposition with $e(R_l)=i$, $e(R_r)=j$.
\paragraph{Lemma 6.17.} Fix $R,c\le|R|$. For each $(I,J;T)$ with $|V(T)\cup I\cup J|\le\tau$ and $R\subseteq I\cap J$, there is exactly one outer-canonical product in
\[\widetilde L_R[D^\tau Q_{c,0}^R D^\tau](\widetilde L_R)^\top\tag{6.24}\]
corresponding to its canonical decomposition, giving the term $M_c^R(I,J;T)\chi_T$.
\begin{proof}Every triple in(6.24) is inner-canonical by the definitions of $\widetilde L_R,Q_{c,0}^R$, hence outer-canonical triples 1-1 correspond to $(I,J;T)$ by canonical decomposition. Fix such a decomposition. $(I,A;T')$ appears once in $\widetilde L_R(I,A)$ in block $L_{R,e_l}$, $e_l=e_{I,A}(T')$; symmetrically for $(J,B;T'')$, $e_r=e_{J,B}(T'')$. Its unique outer-canonical product coefficient is
\[L_{R,e_l}(I,A;T')(\omega/n)^{|A|/2}Q_{c,0}^R(A,B;T_m)(\omega/n)^{|B|/2}L_{R,e_r}(J,B;T'').\tag{6.25}\]
By(6.20),(6.23), if $a=|V(T)\cup I\cup J|\le\tau$, this equals $(\omega/n)^aY_c(|R|,a)=M_c^R(I,J;T)$ by(6.13),(6.16), since canonicality gives $a=e_l+|V(T_m)\cup A\cup B|+e_r$.\end{proof}
\paragraph{Definition 6.18 (First error-matrices).} Let $E_{c,1;\mathrm{negl}}$ sum all outer-canonical products in(6.24) exceeding degree, namely $|V(T)\cup I\cup J|>\tau$. Let the non-canonical part sum all non-outer-canonical products. Lemma6.17 says $M_c^R=[\widetilde L_R(D^\tau Q_{c,0}^R D^\tau)(\widetilde L_R)^\top]_{\mathrm{can}}$, equivalently
\begin{align*}M_c^R&=\widetilde L_R(D^\tau Q_{c,0}^R D^\tau)(\widetilde L_R)^\top\\&\quad-[\widetilde L_R(D^\tau Q_{c,0}^R D^\tau)(\widetilde L_R)^\top]_{\mathrm{non\text{-}can}}-E^R_{c,1;\mathrm{deg}}.\tag{6.26}\end{align*}
The crucial fact is that the main error approximately factorizes through the same left/right factors; the next subsection extends recursive factorization to do this.
\section*{Author Section 6.3: exact recursive factorization statement}
\paragraph{Lemma 6.19 (Recursive approximate factorization; exact case).} For any fixed $R\in\binom{[n]}{\le d/2}$ and $0\le c\le|R|$,
\[M_c^R=\widetilde L_R[D^\tau(Q_{c,0}^R-Q_{c,1}^R+\cdots\pm Q_{c,d}^R)D^\tau](\widetilde L_R)^\top+E_c^R,\tag{6.27}\]
where:
\begin{enumerate}\item All $Q_{c,k}^R$ are supported on $\mathcal S^R\times\mathcal S^R$, with $\mathcal S^R$ from(6.22).
\item $Q_{c,0}^R$ is Definition6.15.
\item For $1<k\le d/2$, within that minor the $(i,j)$ block is
\[Q_{c,k}^R((A,i),(B,j))=\sum_{T_m:|V(T_m)\cup A\cup B|\le\tau}q_{c,k}^R(R_m,i,j)\chi_{T_m},\tag{6.28}\]
$R_m=(A,B;T_m)$, the coefficients symmetric w.r.t. shapes, and
\[\forall(i,j)\quad|q_{c,k}^R(R_m,i,j)|\le\tau^{5\tau}(\omega/n^{1-\epsilon})^{s-p+k/3},\tag{6.29}\]
where $s=(|A|+|B|)/2$ and $p$ is the maximum number of vertex-disjoint paths from $A$ to $B$ in $R_m$.
\item For any $G$, $E_c^R(G)$ is supported on clique rows/columns containing $R$, and w.p. $>1-n^{-9\log n}$,
\[\|E_c^R\|<n^{-\epsilon\tau/2}.\tag{6.30}\]
\end{enumerate}
The proof proceeds via one-round Lemma6.20, followed by its iteration; the complete author proof is continued in the next fragment.
